\section{Correction y el ciclo iterativo}

\subsection{El paso de Correction}

Cuando el Verifier determina que la solution no es válida, se extrae el bug report y se envía al paso de Correction:
\begin{itemize}
    \item Entrada: el problema original + la solution actual + el bug report
    \item Prompt: ``Below is a bug report. If you agree with it, fix the issues. Your output should follow the system prompt instructions.''
    \item Salida: la nueva solution corregida
\end{itemize}

\subsection{El proceso iterativo completo}

Tomando el primer problema como ejemplo:
\begin{enumerate}
    \item Exploración inicial: Step 1 + Step 2 generan la solución inicial
    \item Verificación: el Verifier determina No (hay un bug)
    \item Corrección: se genera una nueva solución a partir del bug report $\rightarrow$ verificación No $\rightarrow$ corrección $\rightarrow$ verificación No
    \item Tras la cuarta corrección se pasa la verificación: correct count = 1
    \item Se pasa la verificación 5 veces consecutivas y finalmente se acepta
    \item En total se pasa por 8 iteraciones, y el ciclo externo solo tiene 1 run
\end{enumerate}

\subsection{Resumen de la sección}

Todo el ciclo iterativo muestra la fuerza del patrón Actor-Critic: aunque la solución inicial tenga varios errores, mediante el ciclo repetido de verificación y corrección finalmente se puede converger a la solución correcta.

